problem:  
Let $\pi: E \to M$ be a smooth $S^k$–bundle over a compact simply connected manifold $M$, with structure group $\mathrm{SO}(k+1)$.  
For each $\varepsilon>0$, let $A_\varepsilon$ be a smooth $\mathrm{SO}(k+1)$–connection on the associated principal bundle, and let $g_\varepsilon$ denote the corresponding **connection metric** on $E$ where vertical directions (tangent to the fibers) are scaled by $\varepsilon^2$ relative to the base metric $h$ on $M$.  

Assume that:  

1. $(M,h)$ admits a Riemannian metric of strictly positive sectional curvature;  
2. for each $\varepsilon$, the connection $A_\varepsilon$ has holonomy contained in a proper subgroup $H < \mathrm{SO}(k+1)$ whose action on the fiber $S^k$ is not transitive (so the associated horizontal distribution is not fat in the sense of Weinstein).  

Question:  
Is there an $\varepsilon_0>0$ and a smooth family of connections $A_\varepsilon$ (possibly varying with $\varepsilon$) such that $(E,g_\varepsilon)$ has **strictly positive sectional curvature** for all sufficiently small $\varepsilon<\varepsilon_0$?  

Equivalently, in the regime of collapsing fiber metrics, can one *restore* positive curvature on the total space by perturbing the connection among non‑fat ones, or is the fatness condition genuinely indispensable for producing positively curved sphere‑bundle metrics?  

Why is it a "good" problem:  
This problem lies at the delicate frontier of geometric bundle constructions. Fatness of the horizontal distribution ensures positivity in O’Neill’s curvature formulas, but whether *approximate* or *variable* fatness can compensate during fiber collapse remains completely unresolved. The question asks whether geometric limits and small connection deformations can bypass the classical fatness obstruction—an issue no current technique fully captures.  

The problem is profound because it forces the analyst to reconsider how horizontal–vertical curvature couplings behave under simultaneous scaling and connection variation. It tightly connects analytic curvature estimates with topological holonomy constraints, bridging differential geometry, global analysis, and bundle topology.  

Moreover, its formulation is simple and transparent: *Can small fiber bundles over positively curved bases acquire positive total curvature without full fatness, provided one can tune the connection?* Any resolution would reshape our understanding of positive curvature in Riemannian submersions—either by revealing a new geometric mechanism or establishing an essential limitation of the fatness condition. This blend of simplicity, depth, and interdisciplinary reach makes it an exemplary “good” mathematical problem.